\section{Zadanie 14}

Zaprosiłeś na przyjęcie 64 gości. Zakładając, że każdy z nich będzie chciał zjeść niezależnie od pozostałych \(0\), \(1\) albo \(2\) kanapki z prawdopodobieństwami \(\frac{1}{4}, \frac{1}{2}\) oraz \(\frac{1}{4}\) odpowiednio, ile kanapek musisz przygotować, aby mieć \(95\%\) pewności, że ich nie zabraknie?

\subsection*{Rozwiązanie}

Niech \(X_i\) oznacza liczbę kanapek zjedzonych przez \(i\)-tego gościa.
\[
	P(X_i=0)=\tfrac14,\quad P(X_i=1)=\tfrac12,\quad P(X_i=2)=\tfrac14
\]

\[
	\ev{X_i}=1,\qquad \variance{X_i}=\tfrac12
\]

Niech \(S=\sum_{i=1}^{64} X_i\). Wówczas
\[
	\ev{S}=64,\qquad \variance{S}=32,\qquad \sigma_S=\sqrt{32}
\]

Z centralnego twierdzenia granicznego:
\[
	S \approx \mathcal N(64,32)
\]

Szukamy \(s\) takiego, że \(P(S\le s)=0.95\). Mamy
\[
	P(S\le s)=P\!\left(\frac{S-64}{\sqrt{32}}\le \frac{s-64}{\sqrt{32}}\right)
\]

Niech \(Z=\frac{S-64}{\sqrt{32}}\sim\mathcal N(0,1)\). Z tablic:
\[
	P(Z\le 1.645)=0.95
\]

Stąd
\[
	\frac{s-64}{\sqrt{32}}=1.645
	\quad\Rightarrow\quad
	s=64+1.645\sqrt{32}\approx 73.3
\]

Zatem należy przygotować co najmniej \(\boxed{74}\) kanapki.
